% =========================================================================
%  Figura: influencia total. Un tubo delgado de líneas de campo une el
%  conductor j (phi_j = 1) con el conductor i (a tierra). La superficie
%  gaussiana cerrada S = S_lat ∪ B_1 ∪ B_2 (a trazos) está formada por la
%  pared del tubo (E tangente, E ⊥ dS) y por dos tapas dentro de los
%  conductores (E = 0). Estilo minimalista: bordes marcados e interiores con
%  un velo radial sutil (estilo «conductorfino»).
% =========================================================================
\begin{figure}[t]
\centering
\begin{tikzpicture}[font=\small, x=1cm, y=1cm]
  % --- conductores --------------------------------------------------------------
  \path[conductorfino]
    (-3.7,1.35) .. controls (-2.6,1.65) and (-1.6,1.2) .. (-1.55,0.3)
    .. controls (-1.5,-0.6) and (-2.4,-1.45) .. (-3.5,-1.25) .. controls (-4.35,-1.05) and (-4.45,1.05) .. cycle;
  \path[conductorfino]
    (2.0,0.95) .. controls (2.6,1.55) and (3.85,1.25) .. (3.95,0.2)
    .. controls (4.05,-0.85) and (3.0,-1.35) .. (2.2,-0.95) .. controls (1.6,-0.65) and (1.6,0.5) .. cycle;
  \node[font=\footnotesize] at (-3.2,0.75) {$i$};
  \node[font=\footnotesize] at (3.3,0.75) {$j$};
  \node[font=\footnotesize, anchor=west] at (3.35,1.45) {$\phi_j=1$};
  \draw[azulacero, line width=0.6pt] (-3.3,-1.27) -- (-3.3,-1.6);
  \pic at (-3.3,-1.6) {tierra};
  \node[font=\footnotesize, anchor=west] at (-3.0,-1.82) {$\phi_i=0$};
  % --- tubo: interior con degradado tenue ---------------------------------------
  \fill[top color=azulvelo!55, bottom color=azulvelo!55, middle color=white]
    (-1.86,0.28) .. controls (-0.6,0.36) and (0.7,0.36) .. (2.0,0.28)
    -- (2.0,-0.28) .. controls (0.7,-0.36) and (-0.6,-0.36) .. (-1.86,-0.28) -- cycle;
  % líneas de campo (de j a i)
  \draw[lcampo] (1.64,0) -- (-1.53,0);
  \draw[azulacero!70, line width=0.4pt] (1.64,0.16) .. controls (0.6,0.2) and (-0.5,0.2) .. (-1.53,0.16);
  \draw[azulacero!70, line width=0.4pt] (1.64,-0.16) .. controls (0.6,-0.2) and (-0.5,-0.2) .. (-1.53,-0.16);
  % --- superficie gaussiana: pared lateral y tapas --------------------------------
  \draw[gauss, line width=0.6pt] (-1.86,0.28) .. controls (-0.6,0.36) and (0.7,0.36) .. (2.0,0.28);
  \draw[gauss, line width=0.6pt] (-1.86,-0.28) .. controls (-0.6,-0.36) and (0.7,-0.36) .. (2.0,-0.28);
  \draw[gauss, line width=0.6pt] (-1.86,0) ellipse (0.07 and 0.28);
  \draw[gauss, line width=0.6pt] (2.0,0) ellipse (0.07 and 0.28);
  % sección transversal del tubo (da la sensación de cilindro)
  \draw[cargapos!60, line width=0.4pt] (0.05,0.34) arc (90:270:0.07 and 0.34);
  \draw[cargapos!60, line width=0.4pt, densely dotted] (0.05,-0.34) arc (-90:90:0.07 and 0.34);
  % --- cargas en los extremos -------------------------------------------------
  \node[signoneg, font=\tiny\bfseries] at (-1.68,0.12) {$-$};
  \node[signoneg, font=\tiny\bfseries] at (-1.68,-0.12) {$-$};
  \node[signopos, font=\tiny\bfseries] at (1.8,0.12) {$+$};
  \node[signopos, font=\tiny\bfseries] at (1.8,-0.12) {$+$};
  % --- rótulos de la superficie --------------------------------------------------
  \draw[grismedio, line width=0.35pt] (-1.86,-0.3) -- (-2.25,-0.75);
  \node[fondo, font=\footnotesize, text=cargapos!85!black] at (-2.4,-0.88) {$B_1$};
  \draw[grismedio, line width=0.35pt] (2.0,-0.3) -- (2.4,-0.75);
  \node[fondo, font=\footnotesize, text=cargapos!85!black] at (2.55,-0.88) {$B_2$};
  \draw[grismedio, line width=0.35pt] (0.9,-0.33) -- (1.15,-0.75);
  \node[fondo, font=\footnotesize, text=cargapos!85!black] at (1.25,-0.9) {$S_{\mathrm{lat}}$};
  \node[font=\scriptsize, text=azulprof] at (-2.75,0.15) {$\vE=\vc{0}$};
  \node[font=\scriptsize, text=azulprof] at (3.05,0.15) {$\vE=\vc{0}$};
  % --- E y dS en un punto de la pared --------------------------------------------
  \coordinate (P) at (-0.5,0.34);
  \draw[-{Stealth[length=3.5pt]}, verdeazul, line width=0.7pt] (P) -- ++(0,0.55);
  \node[fondo, font=\footnotesize, text=verdeazul, anchor=west] at ($(P)+(0.04,0.45)$) {$\dS$};
  \draw[-{Stealth[length=3.5pt]}, azulprof, line width=0.7pt] (P) -- ++(-0.55,0);
  \node[fondo, font=\footnotesize, text=azulprof, anchor=south] at ($(P)+(-0.42,0.04)$) {$\vE$};
  \fill[azulprof] (P) circle (0.9pt);
  % --- leyenda --------------------------------------------------------------------
  \draw[gauss, line width=0.6pt] (-1.3,-1.95) -- (-0.7,-1.95);
  \node[font=\scriptsize, anchor=west] at (-0.6,-1.95) {superficie gaussiana $S=S_{\mathrm{lat}}\cup B_1\cup B_2$};
\end{tikzpicture}
\caption{Influencia total. La superficie cerrada $S$ formada por un tubo de líneas de campo (pared lateral $S_{\mathrm{lat}}$, donde $\vE\perp\dS$) y por
dos tapas $B_1$ y $B_2$ dentro de los conductores (donde $\vE=\vc{0}$) tiene flujo nulo: la carga positiva de la que nace el tubo en $j$ es igual y
opuesta a la negativa en la que muere en $i$.}
\label{fig:tubo}
\end{figure}
